\subsection{Fully lattice-ordered spaces}

DEFINITION 2. --- A Riesz space E is said to be fully lattice-ordered if every nonempty subset of E that is bounded above has a supremum in E.

It is immediate that in a fully lattice-ordered space E, every nonempty subset that is bounded below has an infimum in E.

Examples. --- 1) If A is any set, the space \(R^{A}\) of real-valued functions defined on A is fully lattice-ordered, the supremum in \(R^{A}\) of a family that is bounded above being its upper envelope (GT, IV, §5, No.~5).

\begin{enumerate}
\def\labelenumi{\arabic{enumi})}
\setcounter{enumi}{1}
\tightlist
\item
  Let F be any set; the space \(\mathcal{B}(\mathrm{F})\) of bounded real-valued functions on F, equipped with the order structure induced by that of \(R^{F}\) , is fully lattice-ordered. However, if F is a topological space, the space \(\mathcal{C}(\mathrm{F})\) of continuous real-valued functions on F (equipped with the order structure induced by that of \(R^{F}\) ) is a Riesz space that is not in general fully lattice-ordered (cf.~Exer. 13). Consider for example the case that \(F = R\) ; let I be the interval {]}0,1{[}, \(\varphi_{I}\) the characteristic function of I, and let H be the set of continuous functions \(x(t)\) such that \(x \leqslant \varphi_{I}\) ; it is clear that H is bounded above in \(\mathcal{C}(\mathrm{F})\) . The function \(\varphi_{I}\) is the upper envelope of the \(x \in H\) , but it is not their supremum in \(\mathcal{C}(\mathrm{F})\) , since \(\varphi_{I}\) is lower semi-continuous but not continuous. Let us show that, in fact, H has no supremum in \(\mathcal{C}(\mathrm{F})\) ; it suffices to prove that if u is a continuous function such that \(u \geqslant \varphi_{I}\) , then there exists a continuous function \(v \neq u\) such that \(u \geqslant v \geqslant \varphi_{I}\) . Now, \(u(0) \geqslant 1\) , therefore there exists a number \(\alpha > 0\) such that \(u(t) > 0\) for \(-\alpha \leqslant t \leqslant 0\) ; if w is a continuous function that is zero outside of the interval {]}- \(\alpha,0[\) and is such that \(0 < w(t) < u(t)\) on this interval, then the function v = u - w meets the requirements.
\end{enumerate}

PROPOSITION 1. --- For an ordered vector space E to be fully lattice-ordered, it is necessary and sufficient that E be a Riesz space and that it satisfy one of the following two conditions:

\begin{enumerate}
\def\labelenumi{\alph{enumi})}
\item
  every nonempty subset A, consisting of elements \(\geqslant0\) of E, that is bounded above and directed for the relation \(\leqslant\) , has a supremum in E;
\item
  every nonempty subset A, consisting of elements \(\geqslant 0\) of E and directed for the relation \(\geqslant\) , has an infimum in E.
\end{enumerate}

The conditions are obviously necessary. Conversely, suppose that E is a Riesz space satisfying the condition a). Let B be a nonempty subset of E that is bounded above; the set C consisting of the suprema of the finite subsets of B is directed for the relation \(\leqslant\) ; let a be one its elements and \(C_{a}\) the set of \(x \in C\) that are \(\geqslant a\) ; if we prove that \(C_{a}\) has a supremum then it will also be the supremum of B. Now, \(C_{a} - a\) is a set of elements \(\geqslant 0\) , bounded above and directed for the relation \(\leqslant\) ; it therefore has a supremum b, consequently \(a + b\) is the supremum of \(C_{a}\) .

On the other hand, the condition \(b)\) implies \(a)\) : for, if \(F\) is a nonempty set of elements \(\geqslant 0\) of \(E\) , bounded above and directed for \(\leqslant\) , and if \(c\) is an upper bound for \(F\) , then \(c - F\) is a set of elements \(\geqslant 0\) that is directed for \(\geqslant\) ; if it has an infimum \(m\) , then \(c - m\) is the supremum of \(F\) .

PROPOSITION 2. --- Let E be a Riesz space equipped with a Hausdorff topology that is compatible with its ordered vector space structure (TVS, II, §2, No.~7). If, for every set H ⊂ E that is bounded above and directed for the relation ≤, the section filter of H is convergent, then E is fully lattice-ordered.

Indeed, one knows that the limit of the section filter of H is the supremum of H in E (TVS, II, §2, No.~7, Prop. 18).

